% Part: computability
% Chapter: computability-theory
% Section: prop-reduce

\documentclass[../../../include/open-logic-section]{subfiles}

\begin{document}

\olfileid{cmp}{thy}{ppr}
\olsection{తగ్గించదగినత ధర్మాలు}

$A$ అనేది~$B$కు తగ్గితే $A\leq_m B$ అని రాస్తాం. $A\leq_m B$
అయితే $A$ అనేది~$B$ ``కంటే కష్టం కాదు,'' $B$ అనేది~$A$ ``అంత
కష్టమైనది లేదా మరింత కష్టమైనది'' అని ఈ సంకేతం సూచిస్తుంది. అలాగే
సంఖ్యలపై సాధారణ $\le$ వలెనే $\le_m$, సమితులను (లేదా నిర్ణయ
సమస్యలను) వాటి సంక్లిష్టత ప్రకారం క్రమపరుస్తుందని సూచిస్తుంది. కింది
రెండు ప్రతిపాదనలు ఈ అంతర్బోధకు మద్దతిస్తాయి. మొదటిది $\le_m$
సంక్రమణశీలమని చెబుతుంది.
\sourcecorrection{OLTECOMTHY-013}{ఈ వాక్యంలో ఆంగ్ల మూలం ``says that''ను ``say taht''గా రాసింది. ప్రతిపాదన తెలిపే సంక్రమణశీలతను యథార్థంగా అనువదించాం.}

\begin{prop}
\ollabel{prop:trans-red}
$A\leq_m B$, $B\leq_m C$ అయితే $A\leq_m C$.
\end{prop}

\begin{proof}
$A$ను~$B$కు తగ్గించే ప్రమేయాన్ని, $B$ను~$C$కు తగ్గించే ప్రమేయంతో
సంయుక్తం చేస్తే $A$ను~$C$కు తగ్గించే ప్రమేయం వస్తుంది.
\end{proof}

\begin{prob}
$f$,~$g$లు వరుసగా $A$ నుంచి~$B$కు, $B$ నుంచి~$C$కు అనేకం-ఒకటి
తగ్గింపులైతే, $\comp{f}{g}$ అనేది $A$ నుంచి~$C$కు అనేకం-ఒకటి
తగ్గింపు అని చూపి \olref{prop:trans-red}ను నిరూపించండి.
\end{prob}

\begin{prop}
\ollabel{prop:reduce}
$A$, $B$లు ఏ సమితులైనా, $A\leq_m B$ అని అనుకుందాం.
\begin{enumerate}
\item $B$ గణనీయంగా లెక్కించదగినదైతే $A$ కూడా అలాగే ఉంటుంది.
\item $B$ గణనీయమైతే $A$ కూడా గణనీయం.
\end{enumerate}
\end{prop}

\begin{proof}
$f$ అనేది $A$ నుంచి~$B$కు అనేకం-ఒకటి తగ్గింపు అనుకుందాం. మొదటి
వాదన కోసం, $B$ పాక్షిక ప్రమేయం~$g$ యొక్క నిర్వచన క్షేత్రమైతే,
$A$ అనేది~$\comp{f}{g}$ యొక్క నిర్వచన క్షేత్రమని గమనిస్తే చాలు:
\begin{align*}
x \in A & \text{ అయితేనే } f(x) \in B \\
& \text{ అయితేనే } g(f(x)) \fdefined.
\end{align*}

రెండవ వాదన కోసం, $B$ గణనీయమైతే $B$,~$\Complement{B}$ రెండూ
గణనీయంగా లెక్కించదగినవని గుర్తుచేసుకోండి
(\olref[cmp]{thm:ce-comp}). $f$ అనేది $\Complement{A}$ నుంచి
$\Complement{B}$కు కూడా అనేకం-ఒకటి తగ్గింపని తనిఖీ చేయడం కష్టం కాదు.
కాబట్టి ఈ నిరూపణ మొదటి భాగం ప్రకారం $A$,~$\Complement{A}$ రెండూ
\tetoken{గణనీయంగా లెక్కించదగినవి}{!!{computably
enumerable}}.
అందువల్ల \olref[cmp]{thm:ce-comp} ప్రకారం
$A$ కూడా గణనీయం. (ప్రత్యామ్నాయంగా $\Char{A}=\comp{f}{\Char{B}}$
అని తనిఖీ చేయవచ్చు; $\Char{B}$ గణనీయమైతే $\Char{A}$ కూడా గణనీయం.)
\end{proof}

\begin{prob}
$f$ అనేది $A$ నుంచి~$B$కు అనేకం-ఒకటి తగ్గింపు అనుకుందాం.
$f$ అనేది $\Complement{A}$ నుంచి $\Complement{B}$కు కూడా
అనేకం-ఒకటి తగ్గింపు అని చూపండి.
\end{prob}

\begin{prob}
$f\colon\Nat\to\Nat$ అనేది $A$ నుంచి~$B$కు అనేకం-ఒకటి తగ్గింపు
అయితే, $\Char{A}=\comp{f}{\Char{B}}$ అని చూపండి.
\sourcecorrection{OLTECOMTHY-014}{అనేకం-ఒకటి తగ్గింపును ముందరి నిర్వచనం $\Nat\to\Nat$ ప్రమేయంగా నిర్ణయించినా, ఈ సమస్యలో ఆంగ్ల మూలం దాన్ని $f\colon A\to B$గా తప్పుగా టైప్ చేసింది. నిర్వచనానికీ లక్షణ ప్రమేయాల సంయుక్తానికీ అనుగుణంగా $f\colon\Nat\to\Nat$గా సరిచేశాం.}
\end{prob}

\begin{digress}
ఇతర సందర్భాల్లో, ముఖ్యంగా అనిర్ణయనీయతా ఫలితాలను నిరూపించడానికి,
\emph{ట్యూరింగ్ తగ్గించదగినత} అనే మరింత సాధారణ భావన ఉపయోగపడుతుంది.
\olref[cmp]{cor:comp-k} ప్రకారం $K_0$ పూరకం గణనీయంగా లెక్కించదగినది
కాదు; కనుక అది $K_0$కు అనేకం-ఒకటిగా తగ్గించదగినది కాదని గమనించండి.
అయితే $K_0$ గురించిన ప్రశ్నలకు సమాధానాలు మీకు తెలిస్తే, దాని పూరకం
గురించిన ప్రశ్నల సమాధానాలు కూడా తెలుస్తాయని అంతర్బోధాత్మకంగా స్పష్టం.
$B$ గురించి ప్రశ్నలు అడగగల గణనీయ ప్రక్రియను ఉపయోగించి $A$లోని
ప్రశ్నలకు సమాధానాలు నిర్ణయించగలిగితే, $A$ను~$B$కు ట్యూరింగ్
తగ్గించదగినది అంటాం. ఇది అనేకం-ఒకటి తగ్గించదగినత కంటే ఉదారం;
అనేకం-ఒకటి తగ్గించదగినతలో (1)~$B$ గురించి ఒక్క ప్రశ్ననే అడగాలి,
(2)~``అవును'' అనే సమాధానం $A$ గురించిన ప్రశ్నకు ``అవును'' అనే
సమాధానంగానే మారాలి; ``కాదు''కూ అలాగే జరగాలి. అయినప్పటికీ $A$ అనేది
$B$కు ట్యూరింగ్ తగ్గించదగినదై, $B$ గణనీయమైతే $A$ కూడా గణనీయమే.
(కానీ మనం చూసినట్లుగా, గణనీయంగా లెక్కించదగినతకు సదృశ ప్రకటన నిజం
కాదు.)

మనం చర్చించిన వివిధ తగ్గించదగినతా భావనల గురించి ఆలోచించి, వాటి మధ్య
తేడాలను అర్థం చేసుకోవాలి. అయితే ఈ అధ్యాయంలో మనం అనేకం-ఒకటి
తగ్గించదగినతతోనే వ్యవహరిస్తాం. యాదృచ్ఛికంగా, గత పేరాలో చర్చించిన
రెండు రకాల తగ్గించదగినతలకు గణనా సంక్లిష్టతలో సదృశ భావనలు ఉన్నాయి;
ట్యూరింగ్ యంత్రాలు బహుపది కాలంలో నడవాలనే అదనపు ఆవశ్యకత వాటికి ఉంటుంది.
అనేకం-ఒకటి తగ్గించదగినత యొక్క సంక్లిష్టతా రూపాన్ని \emph{కార్ప్
తగ్గించదగినత} అంటారు; ట్యూరింగ్ తగ్గించదగినత యొక్క సంక్లిష్టతా
రూపాన్ని \emph{కుక్ తగ్గించదగినత} అంటారు.
\end{digress}

\end{document}
